\clearpage
\section{Conclusion}
\label{sec:conclusion}

This thesis studies the labour-market consequences of Brazil's \acrlong{prouni} using a heterogeneity-robust staggered difference-in-differences design applied to administrative employer--employee data for 558 microregions over the period 2005--2019. By combining variation in treatment timing with variation in local scholarship intensity, the analysis estimates reduced-form effects of ProUni exposure on formal wages and employment. The central empirical contribution is not a precise estimate of the individual return to a scholarship, but evidence on whether---and how---local labour markets respond differently to moderate and heavy programme exposure.

Four main findings emerge. First, in Low-Dose microregions the estimated effect on formal-sector wages is positive and economically meaningful: approximately 1.28\% in the baseline specification, rising to 2.59\% under doubly robust estimation. Second, estimated wage effects are weaker in High-Dose microregions and often indistinguishable from zero---a pattern consistent with attenuation of gains in heavily exposed markets, though the high-dose penalty is fragile once macroeconomic geography is controlled for, precluding a definitive saturation interpretation. Third, the age-cohort decomposition provides suggestive evidence on mechanisms: in Low-Dose markets, the results point to positive spillovers extending beyond directly exposed cohorts, while in High-Dose markets they are consistent with relative wage compression. Fourth, the heterogeneity analysis reveals an important asymmetry: ProUni's wage gains are more favourable for non-white workers than for white workers in moderately exposed markets---consistent with the programme's affirmative-action design---yet men appear to benefit more clearly than women in the pooled formal-sector estimates, despite women constituting the majority of scholarship recipients.

These four findings collectively point to a core policy insight: ProUni improves local formal-sector wage outcomes, but its returns depend on market context and do not scale proportionally with programme intensity. The evidence is strongest for positive reduced-form effects in moderately exposed labour markets. It is much weaker for strong claims about saturation, displacement, or exact individual returns.

\subsection*{Policy Implications}

The spatial distribution of ProUni scholarships reveals a structural misalignment between the programme's institutional architecture and where its social returns are likely largest. Because ProUni allocates scholarships through existing private HEI infrastructure, concentration in the South and Southeast is endogenous to commercial higher-education market development rather than to social need \cite{McCowan2007}. The Bel\'india framework \cite{bacha1976, Telles2004} crystallises this tension: market-based mechanisms calibrated to existing infrastructure will underprovide access precisely where skilled labour scarcity---and therefore the marginal social return to each additional graduate---is greatest.

If the positive reduced-form effects in Low-Dose microregions are indicative of this scarcity premium, then a spatially redistributive scholarship allocation---explicitly targeting North, Northeast, and Centre-West microregions with high population share and low private HEI density---would raise the aggregate social return to programme spending. Such reallocation would, however, require investments in local absorptive capacity: firm capacity to hire formally, transport and housing infrastructure to retain graduates, and minimum academic quality standards in newly expanded HEIs. Policy design should therefore treat scholarship volume and labour-market absorptive capacity as complementary inputs rather than substitutes.

The gender asymmetry in wage returns is an independent policy concern. ProUni's majority-female scholarship share suggests strong programme take-up by women, yet formal-sector wage gains overwhelmingly accrue to male workers. This divergence reflects downstream barriers---occupational segregation, firm-level discrimination, and the penalties associated with labour-market re-entry after childbearing---that are orthogonal to the credential access channel \cite{BusteloEtAl2021, Goldin2014}. Complementary interventions targeting these barriers, including occupational desegregation policies and formal childcare provision, would be necessary to realise the welfare gains that female scholarship uptake implies.

\subsection*{Limitations}

Several limitations remain. The analysis is restricted to the formal labour market in RAIS and cannot capture informal employment, self-employment, or non-employment, so welfare effects through formalisation are likely understated \cite{EngbomEtAl2022}. Treatment is measured as scholarship grants rather than degree completions, introducing attenuation through dropout that weakens the first-stage link to tertiary attainment \cite{INEP2024}. The weak first stage prevents credible recovery of a Wald-style individual return to a ProUni degree. Finally, the estimates should be read as ProUni's contribution within a broader policy mix: spatially heterogeneous interactions with FIES and the federal quota law cannot be fully separated using the current design.

\subsection*{Future Research}

Three extensions would most productively build on this thesis. Most importantly, linking individual ProUni scholarship records to RAIS employment histories via the CPF identifier would allow direct estimation of beneficiary-level returns and would resolve the aggregation and attenuation problems that limit the current design. This merge is now technically feasible through the \emph{Base dos Dados} infrastructure \cite{Dahis2022, basedosdados.org} and would transform the reduced-form regional evidence into a precisely estimated LATE for scholarship receipt. A second extension would integrate informal-sector information from PNAD Cont\'inua to test whether the primary channel is wage growth, formalisation, or both \cite{IBGE2025NotasTecnicas}. A third extension would jointly study ProUni, FIES, and the federal quota law in a unified policy evaluation framework, disentangling the contribution of each programme to Brazil's educational transformation and clarifying whether the scarcity premium identified here survives once FIES exposure is directly controlled.
